% ---
% id: meat-core/frank-wolfe/thm:firstorder
% kind: thm
% title: "The estimate and the first-order lower bound"
% summary: >-
%   In exact real arithmetic the best linearisation over the iterates, even
%   ones feasible only up to windows, bounds h_low from below given positive
%   probabilities, maximal branch support, unperturbed gradients and bounds on
%   both linear minima. An exactly feasible point's objective bounds it from
%   above.
% user_statement: >-
%   The headline key rate rests on this bound, which needs no exact optimum.
%   The per-round bound h_low is the minimum of a convex optimisation. At each step
%   whose positivity and full-support conditions hold, the tangent plane of
%   the objective lies below it over the feasible set, so its lowest point
%   bounds h_low from below; the engine reports the best such bound, h_lb.
%   That point is bounded without trusting a solver: one part is
%   computed exactly, the other verified in exact arithmetic from the
%   solver's dual against the exact constraint rows, the floating-point
%   frame W_t's error in those rows deducted. The engine evaluates the
%   objective, its gradient and the maps and frame they are built from in
%   floating point, with no error budget, so a computed h_lb can exceed h_low
%   by an amount no claim bounds. Its points are feasible only within a
%   window, so its reported estimate need not lie above h_low.
% prominence: main
% status: implemented
% uses: [lem:fo-khyp, meat-core/def:branch-support, meat-core/lem:branch-support, meat-core/thm:gradient, meat-core/cor:full-gradient, def:relaxed-feasible-set, prop:relaxed-set-properties, meat-core/thm:choi, meat-core/def:feasible, lem:fo-face, meat-core/thm:convexity, meat-core/prop:value-range, meat-core/thm:branch-formula, meat-core/lem:estimate, def:first-order-setting, meat-core/def:declaration, meat-core/lem:joint-sectors, meat-core/def:admissible-declaration, meat-core/def:admissible-acceptance-set]
% mentions: [rem:acc-side, rem:iterate-relaxed-oracle-exact, req:iterate-checks, lem:feasible-side-dual, req:verified-feasible-side]
% assumes: []
% refs: []
% provenance: newly-derived
% implements: certiqs_key.solvers.frank_wolfe.joint_driver:run_joint_frank_wolfe
% lean:
% ---

The \emph{first-order bound} $\hlb$ is the largest, over the iterates the
optimiser passed through, of the value of the linearisation of the
objective at that iterate, minimised over the feasible set. Part~(ii)
proves $\hlb\le\hlow$ in exact real arithmetic, at maximal support, with
no perturbation of the iterate and given the two inequalities of
\eqref{eq:lmo}; part~(iii) admits an iterate that lies only in
the \tref{relaxed feasible set}, which is where a run's iterates lie
(Remark~\ref{rem:iterate-relaxed-oracle-exact}). Part~(i) restates
Lemma~\ref{lem:estimate}\,(ii) of the core, because it is the whole
warrant of the number a run reports as its estimate, and the membership
$x\in\acSet\times\Feas^{J}$ is exactly what must be checked for that
warrant to hold (Requirement~\ref{req:iterate-checks}). That $\acSet$ is an
\tref{admissible acceptance set} (Definition~\ref{def:admissible-acceptance-set}) is
a hypothesis of the theorem itself and not only of the claims it consumes:
convexity of $\acSet$ is used in Lemma~\ref{lem:fo-khyp}\,(i) and
closedness in Step~2 of the proof. Part~(iv) lets a lower bound on either
minimum stand in for the pairing on its side: on the acceptance side a
certified bound (Remark~\ref{rem:acc-side}), on the feasible side the bound
of Lemma~\ref{lem:feasible-side-dual}, formed as
Requirement~\ref{req:verified-feasible-side} says. Where both stand in, the
bracket reads no point an oracle returned, and the theorem asks nothing of
an oracle. Part~(v) makes the bound independent of
condition~(vii) of an \tref{admissible declaration}, the non-empty feasible set,
which the engine does not decide on every declaration: on an empty feasible
set $\hlow$ is $+\infty$, and every real number bounds it.

\begin{theorem}
\label{thm:firstorder}
Let $\alpha\in(1,\infty)$, let $\Decl$ satisfy
Definition~\ref{def:admissible-declaration}, in the \tref{basic format} or in the general
format of Definition~\ref{def:declaration}, and let $\acSet$ satisfy
Definition~\ref{def:admissible-acceptance-set}. In the general format $\Feas^{J}$ is the set of
lists $(J_{(\mathsf t,s)})$ over the \tref[joint sector]{joint sectors} of
Theorem~\ref{thm:choi}, with $J_s\coloneqq\bigoplus_{\mathsf t}
J_{(\mathsf t,s)}$ for every sector $s$ of the measured system, and
$\Feas^{J}(\varepsilon)$ is read as
Definition~\ref{def:relaxed-feasible-set} reads it there.
\begin{enumerate}[label=(\roman*),leftmargin=2.2em]
\item \emph{(The estimate over-estimates; Lemma~\ref{lem:estimate}\,(ii)
of the core, restated.)} For every $x\in\acSet\times\Feas^{J}$,
\begin{equation}\label{eq:fo-estimate}
  O^{J}_{\alpha}(x)\;\ge\;\hlow .
\end{equation}
\item \emph{(The first-order lower bound.)} Let $\Itset$ be a finite
non-empty index set, and let
$x_t=(\mathbf q_t,J_t)\in\acSet\times\Feas^{J}$, $t\in\Itset$, satisfy
for each $t$:
\begin{itemize}[leftmargin=2.6em]
\item[(H-POS)] $q_t(\hat c)>0$ and $P^{J}_{\hat c}(J_t)>0$ for every
  $\hat c\in\Chacc$, and $K^{J}_{\hat c}(J_t)>0$ for every
  $\hat c\in\Chkey\cap\Chacc$, with $\Chacc$ the \tref{acceptance support}
  \eqref{eq:chacc-def};
\item[(H-SUPP)] every key-active branch operator
  $X_{\hat c,s}\coloneqq\cG^{J}_{\hat c,s}(J_{t,s})$,
  $\hat c\in\Chkey$, $s\in\SecSet_K(\hat c)$, attains maximal support:
  $\operatorname{supp}X_{\hat c,s}=W_{\hat c,s}$, the
  \tref{maximal branch support} \eqref{eq:Wcs}, equivalently
  $\operatorname{rank}X_{\hat c,s}=\dim W_{\hat c,s}
  =\operatorname{rank}\mathbf K_{\hat c,s}$ by \eqref{eq:Kstack-rank},
  the inclusion $\subseteq$ holding always by \eqref{eq:Wcs-bound};
  $J_{t,s}\succ0$ suffices for it (Lemma~\ref{lem:branch-support}\,(i));
\item[(H-NOPERT)] the pair
  $(\lmoq_t,\lmoJ_t)\coloneqq\nabla O^{J}_{\alpha}(x_t)$ is the
  split-form gradient of Theorem~\ref{thm:gradient}\,(i)(a)--(b) assembled
  by Corollary~\ref{cor:full-gradient} \emph{at $x_t$ itself}: no
  perturbation of $x_t$, of any $J_{t,s}$, or of any branch operator
  built from them is applied before the gradient is taken.
\end{itemize}
Let $(v_t,S_t)$ satisfy both inequalities of \eqref{eq:lmo} at
$(\lmoq_t,\lmoJ_t)$ (Definition~\ref{def:first-order-setting}): the
\tref{acceptance-side inequality} and the
\tref{feasible-side inequality}.
Then, \emph{in exact real arithmetic},
\begin{equation}\label{eq:hlb}
  \hlb\;\coloneqq\;\max_{t\in\Itset}\Big[
  O^{J}_{\alpha}(x_t)
  +\langle\lmoq_t,v_t-\mathbf q_t\rangle
  +\langle\lmoJ_t,S_t-J_t\rangle\Big]
  \;\le\;\hlow ,
\end{equation}
and $\max\{0,\hlb\}\le\hlow$ as well.
\item \emph{(The iterate read at a tolerance, the minima over the exact
set.)} Let $\varepsilon\ge0$ and let $\Feas^{J}(\varepsilon)$ be the
relaxed feasible set of Definition~\ref{def:relaxed-feasible-set}.
\begin{enumerate}[label=(iii\alph*),leftmargin=2.4em]
\item \emph{The estimate.} For every
$x\in\acSet\times\Feas^{J}(\varepsilon)$,
$O^{J}_{\alpha}(x)\ge\hlow[\Feas^{J}(\varepsilon)]$, the Choi-form
relaxed infimum \eqref{eq:hlow-feasJ}; and
$\hlow[\Feas^{J}(\varepsilon)]\le\hlow\le\hqkd$ by
\eqref{eq:feas-tol-mono-J}.
\item \emph{The bound.} Let $x_t\in\acSet\times\Feas^{J}(\varepsilon)$,
$t\in\Itset$, satisfy (H-POS), (H-SUPP) and (H-NOPERT) as stated in (ii),
and let $(v_t,S_t)$ satisfy \eqref{eq:lmo} \emph{as written}, its
acceptance-side inequality and its feasible-side inequality, with the two
minima over $\acSet$ and over the \emph{exact} $\Feas^{J}$. Then, in exact
real arithmetic,
$\hlb\le\hlow$ and $\max\{0,\hlb\}\le\hlow$, with $\hlb$ of
\eqref{eq:hlb}.
\end{enumerate}
\item \emph{(Lower bounds in place of the pairings.)} The conclusions of
(ii) and of (iiib) hold as well when, at some or all $t\in\Itset$, the
pairing $\langle\lmoq_t,v_t\rangle$ in \eqref{eq:hlb} is replaced by a real
number $b^{\mathrm A}_t$, or the pairing $\langle\lmoJ_t,S_t\rangle$ by a
real number $b^{\mathrm F}_t$, or both, with
\begin{equation}\label{eq:lmo-bounds}
  b^{\mathrm A}_t\;\le\;\min_{\mathbf q'\in\acSet}\langle\lmoq_t,\mathbf q'\rangle
  \qquad\text{and}\qquad
  b^{\mathrm F}_t\;\le\;\min_{J'\in\Feas^{J}}\langle\lmoJ_t,J'\rangle ,
\end{equation}
each inequality asked only where its number is used, the second minimum
taken over the \emph{exact} $\Feas^{J}$, and the pairings
$\langle\lmoq_t,\mathbf q_t\rangle$ and $\langle\lmoJ_t,J_t\rangle$ being
kept. At such a $t$ the inequality of \eqref{eq:lmo} on a replaced side is
not asked of $(v_t,S_t)$, and where both sides are replaced no pair
$(v_t,S_t)$ is asked at all: the bracket of \eqref{eq:hlb} at that $t$ is
then
\begin{equation*}
  O^{J}_{\alpha}(x_t)+b^{\mathrm A}_t-\langle\lmoq_t,\mathbf q_t\rangle
  +b^{\mathrm F}_t-\langle\lmoJ_t,J_t\rangle .
\end{equation*}
\item \emph{(An empty feasible set.)} In place of the hypothesis on $\Decl$
above, let $\Decl$ be a tuple of the form of
Definition~\ref{def:declaration} that satisfies conditions (i)--(vi) and
(viii) of Definition~\ref{def:admissible-declaration}, and form $\Feas$, the
objective $O_{\alpha}$ and $\hlow$ from it by the formulas of
Definition~\ref{def:feasible}. Let $\Feas=\varnothing$. Then $\hlow=+\infty$,
so every real number is a lower bound on $\hlow$, the value of
\eqref{eq:hlb} that a run forms at its iterates included. Consequently the
conclusions $\hlb\le\hlow$ and $\max\{0,\hlb\}\le\hlow$ of (ii), (iiib) and
(iv), for a real $\hlb$, hold of every tuple of the form of
Definition~\ref{def:declaration} that satisfies conditions (i)--(vi) and
(viii) of Definition~\ref{def:admissible-declaration}, whether or not it
satisfies condition~(vii): under the hypotheses of those parts where it does,
and because $\hlow=+\infty$ where it does not.
\end{enumerate}
\end{theorem}

\begin{proof}
\emph{Idea.} At each iterate the objective has a gradient, supplied by
Corollary~\ref{cor:full-gradient} with its positivity weakened to the
acceptance support by Lemma~\ref{lem:fo-khyp}, and at maximal support that
gradient is a genuine directional derivative in every feasible direction
(Lemma~\ref{lem:fo-face}). Convexity of the objective turns the
derivative into a linear minorant; its minimum over the product
$\acSet\times\Feas^{J}$ separates into the two linear minimisations of
\eqref{eq:lmo}; and a maximum of lower bounds is a lower bound.
In (iiib) the point lies in the relaxed set while the comparison point
lies in the exact one, and the chord argument runs in the relaxed set,
which is convex.

\emph{(i).} This is Lemma~\ref{lem:estimate}\,(ii) of the core.

\emph{(ii).} Fix $t$ and write $x\coloneqq x_t$,
$\nabla O\coloneqq(\lmoq_t,\lmoJ_t)$.

\emph{Step 1 (the gradient exists and is a genuine directional derivative
in every feasible direction).} By (H-POS) and (H-SUPP), and by
Lemma~\ref{lem:fo-khyp}\,(iii), which is what lets (H-POS) be assumed
on $\Chacc$ and its $K$-part only on $\Chkey\cap\Chacc$,
Corollary~\ref{cor:full-gradient} applies at $x$ and supplies
$\nabla O$ through \eqref{eq:grad-q}--\eqref{eq:grad-pieces}, with
$\lmoq_t\in\mathbb{R}^{\Chacc}$; by
(H-NOPERT) that is the object the bracket uses. Let
$y=(\mathbf q',J')\in\acSet\times\Feas^{J}$ and put $h\coloneqq y-x$. Its
$J$-component $\Delta_s\coloneqq J'_s-J_{t,s}$ is a difference of positive
semidefinite operators and hence Hermitian, so by
Lemma~\ref{lem:fo-face} every branch operator's direction
$\cG^{J}_{\hat c,s}(\Delta_s)$ is supported inside
$\operatorname{supp}X_{\hat c,s}$, and $t'\mapsto T(X_{\hat c,s}
+t'\cG^{J}_{\hat c,s}(\Delta_s))$ is differentiable at $t'=0$. The
remaining dependence of $O^{J}_{\alpha}$ on $x$ is through the affine maps
$\mathbf q\mapsto q(\hat c)$ and $J\mapsto P^{J}_{\hat c}(J)$ and through
finitely many compositions of $\log_2$ with quantities that (H-POS) keeps
strictly positive, all differentiable there, the sums running over
$\Chacc$, respectively $\Chkey\cap\Chacc$, by
Lemma~\ref{lem:fo-khyp}\,(ii), which deletes the remaining terms
identically on $\acSet\times\Feas^{J}$ rather than bounding them. Hence
$t'\mapsto O^{J}_{\alpha}(x+t'h)$ is differentiable at $t'=0$ with
derivative $\langle\nabla O,h\rangle$, by the chain and sum rules, exactly
as in the proof of Corollary~\ref{cor:full-gradient}, whose
maximal-support step is Lemma~\ref{lem:fo-face} applied to that same
$\Delta_s$. Theorem~\ref{thm:gradient}\,(ii) is what makes this a
\emph{complete} statement rather than an assumption: the only way a
one-sided derivative of $T$ can fail to be finite is a support-enlarging
direction, and Lemma~\ref{lem:fo-face}\,(b) says there is none.

\emph{Step 2 (convexity turns the derivative into an inequality).} By
Theorem~\ref{thm:convexity}, $O^{J}_{\alpha}$ is convex with values in
$(-\infty,+\infty]$ on the convex set $\acSet\times\Feas^{J}$:
Theorem~\ref{thm:convexity} states this for $O_{\alpha}$ on
$\Delta(\Chatset)\times\Feas$, and $O^{J}_{\alpha}=O_{\alpha}\circ
(\idmap\times\chi_{\Qsys})$ with $\chi_{\Qsys}$ \emph{linear}, so the
composition is convex on the convex preimage, a pullback along a linear
map that needs no injectivity, which is Theorem~\ref{thm:choi}'s own
convexity clause. For a convex extended-real function the difference
quotient $t'\mapsto\big(O^{J}_{\alpha}(x+t'h)-O^{J}_{\alpha}(x)\big)/t'$
is non-decreasing on $(0,1]$: this is the three-chord inequality, and it
holds with the value $+\infty$ admitted. $O^{J}_{\alpha}(x)$ is finite by
(H-POS), and $x+t'h=(1-t')x+t'y$ stays in the convex set
$\acSet\times\Feas^{J}$ for $t'\in[0,1]$, so the quotient is defined
there. Its limit as $t'\downarrow0$ is $\langle\nabla O,h\rangle$ by
Step~1, and its value at $t'=1$ is
$O^{J}_{\alpha}(y)-O^{J}_{\alpha}(x)$; monotonicity gives
$\langle\nabla O,h\rangle\le O^{J}_{\alpha}(y)-O^{J}_{\alpha}(x)$, that
is
\begin{equation}\label{eq:fo-subgrad}
  O^{J}_{\alpha}(y)\;\ge\;O^{J}_{\alpha}(x)
  +\langle\nabla O,\,y-x\rangle
  \qquad\text{for every }y\in\acSet\times\Feas^{J},
\end{equation}
including every $y$ at which $O^{J}_{\alpha}=+\infty$, where the
inequality is trivial. Taking the infimum over $y$ and using
Theorem~\ref{thm:choi},
\begin{equation}\label{eq:fo-inf}
  \hlow\;\ge\;O^{J}_{\alpha}(x)
  +\min_{y\in\acSet\times\Feas^{J}}\langle\nabla O,\,y-x\rangle ,
\end{equation}
the minimum being attained because $\acSet\times\Feas^{J}$ is compact and
the map is linear. That $\Feas^{J}$ is compact is read off
\eqref{eq:choi-feas} directly: it is closed, being an intersection of the
positive semidefinite cones with affine equalities and inequalities; and
it is bounded, because $J_s\succeq0$ with
$\sum_s\tr_{\Qsys_s}J_s=\idon{A_0}$ forces
$\sum_s\tr J_s=\tr\idon{A_0}=d_{A_0}$ and hence
$\|J_s\|\le\tr J_s\le d_{A_0}$ for every $s$. $\acSet$ is
closed in the simplex by Definition~\ref{def:admissible-acceptance-set}. In untrusted-source
mode the same reading applies to $\Feas$ with $\tr\omega=1$ in place of
the coupling condition.

\emph{Step 3 (the minimum separates over the product).} The map
$y\mapsto\langle\nabla O,y\rangle
=\langle\lmoq_t,\mathbf q'\rangle+\langle\lmoJ_t,J'\rangle$ is a sum of
one term in each factor, and $\acSet\times\Feas^{J}$ is a product set, so
\begin{equation*}
  \min_{y\in\acSet\times\Feas^{J}}\langle\nabla O,y\rangle
  =\min_{\mathbf q'\in\acSet}\langle\lmoq_t,\mathbf q'\rangle
  +\min_{J'\in\Feas^{J}}\langle\lmoJ_t,J'\rangle
  \;\ge\;\langle\lmoq_t,v_t\rangle+\langle\lmoJ_t,S_t\rangle ,
\end{equation*}
the last step being \eqref{eq:lmo}. The pair $(v_t,S_t)$ is \emph{not}
required to be feasible, only to satisfy the two inequalities.
Substituting into \eqref{eq:fo-inf},
\begin{equation*}
  \hlow\;\ge\;O^{J}_{\alpha}(x_t)
  +\langle\lmoq_t,v_t-\mathbf q_t\rangle
  +\langle\lmoJ_t,S_t-J_t\rangle .
\end{equation*}

\emph{Step 4 (a maximum of lower bounds is a lower bound).} The display
holds for each $t\in\Itset$, which is finite and non-empty, so it holds
for the largest of the right-hand sides, which is \eqref{eq:hlb}. Finally
$\hlow\ge0$ by Proposition~\ref{prop:value-range}\,(i), whose proof
consumes only Definition~\ref{def:admissible-declaration} and
Theorem~\ref{thm:branch-formula}; so $\max\{0,\hlb\}$ is the maximum of
two lower bounds on $\hlow$ and is therefore one.

\emph{(iiia).} $\hlow[\Feas^{J}(\varepsilon)]$ is by \eqref{eq:hlow-feasJ}
the infimum of $O^{J}_{\alpha}$ over $\acSet\times\Feas^{J}(\varepsilon)$
and $x$ is a point of that set, which is the whole of (i)'s argument with
the set substituted; nothing else of (i) is used, and in particular no
property of the set is. That
$\hlow[\Feas^{J}(\varepsilon)]\le\hlow\le\hqkd$ is
\eqref{eq:feas-tol-mono-J}, whose one substantive step is the inclusion
$\Feas^{J}\subseteq\Feas^{J}(\varepsilon)$ proved there.

\emph{(iiib).} The proof is (ii)'s four steps with one change, made in
Step~2: the point $x_t$ ranges over the \emph{relaxed} set while the
comparison point $y$ ranges over the \emph{exact} one. Fix
$t$ and write $x\coloneqq x_t$, $\nabla O\coloneqq(\lmoq_t,\lmoJ_t)$.

\emph{Step 1 is unchanged, and reads only per-block positivity.} By
(H-POS) and (H-SUPP) and Lemma~\ref{lem:fo-khyp}\,(iii),
Corollary~\ref{cor:full-gradient} applies at $x$ and supplies $\nabla O$.
That the corollary applies here, at a point of the relaxed set, is a
property of its hypothesis and not of this step: it asks $\mathbf q\in\acSet$ and
$J_s\in\Pos{A_0\Qsys_s}$ per block, names no normalisation, no coupling
condition and no declared row, and every point of $\Feas^{J}(\varepsilon)$
meets it, per-block positive semidefiniteness being the one condition the
relaxed set does not relax. And Lemma~\ref{lem:fo-khyp}\,(iii), the
statement that weakens the corollary's positivity clause to $\Chacc$, is
itself stated on $\acSet\times\Feas^{J}(\varepsilon)$ for every
$\varepsilon\ge0$, so the weakening reaches this point too.
By (H-NOPERT) the object so supplied is the one the bracket uses. Let
$y=(\mathbf q',J')\in\acSet\times\Feas^{J}$, the \emph{exact} set,
and put $h\coloneqq y-x$. Its $J$-component
$\Delta_s\coloneqq J'_s-J_{t,s}$ is a difference of two positive
semidefinite operators, $J'_s$ from $\Feas^{J}$ and $J_{t,s}$ from
$\Feas^{J}(\varepsilon)$ (per-block positive semidefiniteness being the
one condition the relaxed set does not relax,
Definition~\ref{def:relaxed-feasible-set}), hence Hermitian. That is all
Lemma~\ref{lem:fo-face} asks: its hypotheses are $J_s\in\Pos{A_0\Qsys_s}$
and $\Delta$ Hermitian, and it names no feasible set, no normalisation,
no coupling condition and no declared row. So every branch operator's
direction $\cG^{J}_{\hat c,s}(\Delta_s)$ is supported inside
$\operatorname{supp}X_{\hat c,s}$ and the map
$t'\mapsto T(X_{\hat c,s}+t'\cG^{J}_{\hat c,s}(\Delta_s))$ is
differentiable at $t'=0$. The remaining dependence is through the affine
$\mathbf q\mapsto q(\hat c)$ and $J\mapsto P^{J}_{\hat c}(J)$ and through
compositions of $\log_2$ with quantities (H-POS) keeps strictly positive.
Lemma~\ref{lem:fo-khyp}\,(ii) deletes the terms indexed outside $\Chacc$
identically here as well: its input is $0\le K^{J}_{\hat c}\le
P^{J}_{\hat c}$ on the set the points lie in, which
Proposition~\ref{prop:relaxed-set-properties}\,(c) establishes on the
whole cone (\emph{First}), from positivity alone and consuming no trace
hypothesis, and in the Choi form on $\Feas^{J}(\varepsilon)$. Hence
$t'\mapsto O^{J}_{\alpha}(x+t'h)$ is differentiable at $t'=0$ with
derivative $\langle\nabla O,h\rangle$.

\emph{Step 2 (the segment lies in the relaxed set, and the chord argument
runs there).} $\acSet\times\Feas^{J}(\varepsilon)$ is convex
(Proposition~\ref{prop:relaxed-set-properties}\,(b); $\acSet$ by
Definition~\ref{def:admissible-acceptance-set}), and $O^{J}_{\alpha}$ is convex on it with
values in $(-\infty,+\infty]$ by
Proposition~\ref{prop:relaxed-set-properties}\,(c), whose Choi form pulls
Theorem~\ref{thm:convexity} back along the \emph{linear} $\chi_{\Qsys}$,
needing no injectivity. Both endpoints lie in that convex set:
$x\in\acSet\times\Feas^{J}(\varepsilon)$ by hypothesis, and
$y\in\acSet\times\Feas^{J}\subseteq\acSet\times\Feas^{J}(\varepsilon)$ by
Proposition~\ref{prop:relaxed-set-properties}\,(a). So $x+t'h=(1-t')x+t'y$
stays in it for $t'\in[0,1]$, $O^{J}_{\alpha}(x)$ is finite by (H-POS),
and the three-chord inequality gives, exactly as in (ii),
$\langle\nabla O,h\rangle\le O^{J}_{\alpha}(y)-O^{J}_{\alpha}(x)$, that
is, \eqref{eq:fo-subgrad} \emph{for every $y\in\acSet\times\Feas^{J}$}.
Taking the infimum over that set, and not over the relaxed one, and using
Theorem~\ref{thm:choi}, which says the infimum of $O^{J}_{\alpha}$ over
$\acSet\times\Feas^{J}$ is $\hlow$,
\begin{equation}\label{eq:fo-inf-hybrid}
  \hlow\;\ge\;O^{J}_{\alpha}(x)
  +\min_{y\in\acSet\times\Feas^{J}}\langle\nabla O,\,y-x\rangle ,
\end{equation}
the minimum attained because $\acSet\times\Feas^{J}$ is compact and the
map linear, compactness being read off \eqref{eq:choi-feas} exactly as in
(ii), and non-emptiness following from Definition~\ref{def:admissible-declaration}\,(vii)
with Theorem~\ref{thm:choi} mapping $\Feas^{J}$ onto $\Feas$.

\emph{Steps 3 and 4 are those of (ii).} The map
$y\mapsto\langle\nabla O,y\rangle$ separates over the product
$\acSet\times\Feas^{J}$, and the two minima that appear are exactly the
two of \eqref{eq:lmo} as the hypothesis of (iiib) states them; substituting
and maximising over the finite non-empty $\Itset$ gives $\hlb\le\hlow$.
Finally $\hlow\ge0$ by Proposition~\ref{prop:value-range}\,(i), a
statement about the \emph{exact} set that is not claimed at
$\varepsilon>0$ (Proposition~\ref{prop:relaxed-set-properties}\,(c),
(c$''$)) and is not needed there, so $\max\{0,\hlb\}$ is the maximum of
two lower bounds on $\hlow$ and is one. Part~(ii) is the case
$\varepsilon=0$.

\emph{(iv).} Steps~1, 2 and~4, of (ii) and of (iiib), do not read the pair
$(v_t,S_t)$, and Step~3 reads of it only the two inequalities of
\eqref{eq:lmo}, the first of its acceptance half and the second of its
feasible half, each to bound one of the two minima into which the minimum
over the product separates. At a $t$ where a replacement is made, put
$a_t\coloneqq b^{\mathrm A}_t$ if the acceptance side is replaced and
$a_t\coloneqq\langle\lmoq_t,v_t\rangle$ otherwise, and
$f_t\coloneqq b^{\mathrm F}_t$ if the feasible side is replaced and
$f_t\coloneqq\langle\lmoJ_t,S_t\rangle$ otherwise. The separation of
Step~3, with \eqref{eq:lmo-bounds} on each replaced side and \eqref{eq:lmo}
on each side not replaced, gives
\begin{equation*}
  \min_{y\in\acSet\times\Feas^{J}}\langle\nabla O,y\rangle
  =\min_{\mathbf q'\in\acSet}\langle\lmoq_t,\mathbf q'\rangle
  +\min_{J'\in\Feas^{J}}\langle\lmoJ_t,J'\rangle
  \;\ge\;a_t+f_t ,
\end{equation*}
and substituting into \eqref{eq:fo-inf}, respectively
\eqref{eq:fo-inf-hybrid}, gives
$\hlow\ge O^{J}_{\alpha}(x_t)+a_t-\langle\lmoq_t,\mathbf q_t\rangle
+f_t-\langle\lmoJ_t,J_t\rangle$, which is the bracket of \eqref{eq:hlb}
with the replacements made. Where both sides are replaced, neither $v_t$
nor $S_t$ enters any step. Step~4 then runs unchanged.

\emph{(v).} No claim is used for $\hlow=+\infty$. By \eqref{eq:hlow-def},
$\hlow$ is the infimum of $O_{\alpha}$ over $\mathbf q\in\acSet$ and
$\omega\in\Feas$; with $\Feas=\varnothing$ the inner infimum is over the empty
set and is $+\infty$ for every $\mathbf q$, and so is the outer one, whether
$\acSet$ is empty or not. Every real number is below $+\infty$. For the
consequence: where the tuple satisfies condition~(vii) as well, it
satisfies Definition~\ref{def:admissible-declaration} in full, and with it
the two conditions Definition~\ref{def:declaration} places on its raw-key
map, which are conditions (v) and (viii); it is then a protocol declaration
satisfying the hypothesis of (ii), (iiib) and (iv); where it does not, $\Feas=\varnothing$, so
$\hlow=+\infty$ and the real number $\hlb$ is below it.

\emph{What (iiib) uses of the relaxed set.} Only these: that
$\Feas^{J}\subseteq\Feas^{J}(\varepsilon)$
(Proposition~\ref{prop:relaxed-set-properties}\,(a)); that
$\acSet\times\Feas^{J}(\varepsilon)$ is convex (\emph{ibid.}\,(b)); that
$O^{J}_{\alpha}$ is convex on it (\emph{ibid.}\,(c)); that every member is
per-block positive semidefinite (Definition~\ref{def:relaxed-feasible-set}:
the one condition not relaxed, what Lemma~\ref{lem:fo-face} consumes, and
what the hypothesis of Corollary~\ref{cor:full-gradient} asks of $x$, the
corollary being the statement that supplies $(\lmoq_t,\lmoJ_t)$ and
therefore the one whose scope decides whether \eqref{eq:hlb} names an
object at all); and that $0\le K^{J}_{\hat c}\le P^{J}_{\hat c}$ on it
(\emph{ibid.}\,(c), \emph{First}, which is the cone clause of
Theorem~\ref{thm:branch-formula} and what Lemma~\ref{lem:fo-khyp}\,(ii)
consumes). The list does not say that compactness of the relaxed set is
unused: no step of this proof reads it, the minimum being taken over the
exact set, but the convexity of $O^{J}_{\alpha}$ that Step~2 imports is
Proposition~\ref{prop:relaxed-set-properties}\,(c) \emph{Third}, whose own
derivation consumes the convexity \emph{and} the compactness of (b), the
latter through the continuity of $K_{\hat c}$; so compactness is used one
level down and not in the minimisation. Nor does it say that lower
semicontinuity is unavailable: the chord argument needs convexity alone,
but the fact imported carries both. Of the \emph{exact} set the proof uses
compactness, non-emptiness, and that its infimum is $\hlow$
(Theorem~\ref{thm:choi}). No statement is made about
$\hlow[\Feas(\varepsilon)]$ here, and none is used: the state-form relaxed
set enters this part nowhere.

\emph{The general format.} No step above reads a constraint row, the
number of \tref[measured factor]{measured factors} or the number of Alice
sectors, so the proof applies verbatim with $\Feas^{J}$ the joint-sector
Choi set. Step~1 reads the hypothesis of
Corollary~\ref{cor:full-gradient}, per-block positive semidefiniteness of
the point, which names no feasible set, no normalisation and no declared
row; Lemma~\ref{lem:fo-face}, whose hypotheses are
$J_s\in\Pos{A_0\Qsys_s}$ and $\Delta$ Hermitian and which reads only
\eqref{eq:Wcs-bound}, stated for the admissible Alice subspace in either
format; Theorem~\ref{thm:gradient}\,(ii), a statement about the functional
on the positive cone of a branch space; and
Lemma~\ref{lem:fo-khyp}\,(ii) and~(iii), which read only the acceptance
support, per-block positive semidefiniteness and the cone clause. Step~2
reads convexity of $\acSet\times\Feas^{J}$, convexity of the objective on
it, compactness of $\Feas^{J}$ and the identification of the infimum over
it with $\hlow$, all four supplied by the general-format paragraph of
Theorem~\ref{thm:choi}, including for the \tref{party-local} rows, which are
affine semidefinite conditions there like every other row; that paragraph
cites the convexity of the objective from Theorem~\ref{thm:convexity},
whose inputs, convexity of the feasible set and concavity of $K_{\hat c}$,
Definition~\ref{def:feasible} and Lemma~\ref{lem:joint-sectors} supply in
the general format. Steps~3 and~4 read that the variable is a product and
that the gradient pairs term by term over its factors, that $\Itset$ is
finite and non-empty, and, for the clamp at zero,
Proposition~\ref{prop:value-range}\,(i), which consumes only the
conventions of the kernel, the objective and the cone clause. Part~(iiib)
reads in addition (a), (b) and (c) of
Proposition~\ref{prop:relaxed-set-properties} and per-block positive
semidefiniteness of the relaxed set, which hold in the general format by
that proposition's (e); and (iv) reads nothing more.
\end{proof}
